\documentclass[11pts, a4paper]{ctexart}
\usepackage{amsmath}
\usepackage{booktabs}
\usepackage{geometry}
\usepackage{xcolor}
\usepackage{tcolorbox} % 加框强调
\usepackage{tablefootnote}

\geometry{left=2cm, right=2cm, top=2cm, bottom=2cm}

\usepackage{titlesec}
\titleformat{\section}[hang]{\Large\bfseries\raggedright}{\thesection}{1em}{}

\renewcommand{\arraystretch}{1.3}

\begin{document}

\section{导数}

导数有几个不同的记号，$f'(x)$、$\frac{\mathrm{d}}{\mathrm{d}x} f(x)$、$\mathbf D f(x)$ 都可以用来代表求导。如果需要反复求导（高阶导数），那么可以用 $f'''(x)$、$f^{(k)}(x)$、$\frac{\mathrm{d^k}}{\mathrm{d}x^k} f(x)$、$\mathbf D^k f(x)$ 这些记号。

我们做的所有操作都是\textbf{形式上的}。你只需要（熟练）记住以下这些事实：

\subsection*{基本规则}

\begin{center}
\begin{tabular}{lll}
\toprule
规则名称 & 公式 & 备注 \\
\midrule
线性运算的导数 & $(af(x)+bg(x))'=af'(x)+bg'(x)$ & 逐项求导 \\
乘积的导数 & $(f(x)g(x))'=f'(x)g(x)+f(x)g'(x)$ &  每项导一次\\
复合的导数 & $(f(g(x)))'=f'(g(x))g'(x)$ & 链式法则 \\
初等函数的导数 & $(x^{\alpha})'=\alpha x^{\alpha-1}$ & 幂函数求导 \\
 & $(\ln x)'=x^{-1}$ & 对数函数求导 \\
 & $(e^x)'=e^x$ & 指数函数求导 \\
\bottomrule
\end{tabular}
\end{center}

\subsection*{常用结论（练习）}

\begin{center}
\begin{tabular}{lll}
\toprule
名称 & 公式 & 备注 \\
\midrule
常数的导数 & $\frac{\mathrm{d}}{\mathrm{d}x} c=0$ &  \\
商的导数 & $(f(x)/g(x))'=\frac{f'(x)g(x)-f(x)g'(x)}{g(x)^2}$ &  \\
对数复合 & $(\log f(x))'=\frac{f'(x)}{f(x)}$ &  \\
指数复合 & $(e^{f(x)})'=f'(x)e^{f(x)}$ &  \\
乘积的高阶导数 & $(fg)^{(k)}=\sum_i \binom{k}{i}f^{(i)}g^{(k-i)}$ & \\
多项乘积的导数 & $(f_1f_2\cdots f_k)'=(f_1f_2\cdots f_k)\sum_i\frac{f_i'}{f_i}$ & \\
\bottomrule
\end{tabular}
\end{center}

\section{级数运算与数列运算的关系}

\begin{center}
\begin{tabular}{lll}
\toprule
生成函数操作 & 数列对应 & 描述 \\
\midrule
$H(x)=aF(x)+bG(x)$ & $h_n=af_n+bg_n$ & 线性运算 \\
$H(x)=F(x)G(x)$ & $h_n=\sum_{i=0}^n f_ig_{n-i}$ & 卷积 \\
$H(x)=F_1(x)F_2(x)\cdots F_m(x)$ & $h_n=\sum_{i_1+\cdots+i_m=n}f_{i_1}\cdots f_{i_m}$ & 多项卷积 \\
$H(x)=F'(x)$ & $h_n=(n+1)f_{n+1}$ & 求导 \\
$H(x)=x^mF(x)$ & $h_n=f_{n-m}$ & 移位（小心负的 $m$）\\
$H(x)=\frac{1}{1-x}F(x)$ & $h_n=\sum_{i\le n}f_i$ & 前缀和 \\
$H(x)=(1-x)F(x)$ & $h_n=f_n-f_{n-1}$ & 差分 \\
$H(x)=xF'(x)$ & $h_n=nf_n$ & 原地乘 $n$ \\
$H(x)=F(x)\bmod (x^k-1)$ & $h_r=\sum_n f_{kn+r}$ & 同余求和 \\
$H(x)=x^kF^{(k)}(x)$ & $h_n=n^{\underline{k}}f_n$ & 乘下降幂 \\
$H(x)=(x\mathbf D)^kF(x)$ & $h_n=n^k f_n$ & 乘普通幂 \\
\bottomrule
\end{tabular}
\end{center}



\section{最重要的一些生成函数}

\subsection*{基本恒等式}

\begin{center}
\begin{tabular}{lll}
\toprule
函数 & 展开式 & 描述 \\
\midrule
$(1+x)^{\alpha}$ & $\sum_{n\ge 0}\binom{\alpha}{n}x^n$ & 二项式定理 \\
$\ln\left(\frac{1}{1-x}\right)$ & $\sum_{n\ge 1}\frac{x^n}{n}$ & 对数函数 \\
$e^x$ & $\sum_{n\ge 0}\frac{x^n}{n!}$ & 指数函数 \\
\bottomrule
\end{tabular}
\end{center}

\subsection*{常见特殊形式}

\begin{center}
\begin{tabular}{lll}
\toprule
函数 & 展开式 & 描述 \\
\midrule
$\frac{1}{1-x}$ & $\sum_{n\ge 0}x^n$ & 全 $1$ 数列 \\
$\frac{1}{1+x}$ & $\sum_{n\ge 0}(-1)^nx^n$ & 交错数列 \\
$\frac{1}{1-cx}$ & $\sum_{n\ge 0}c^nx^n$ & 等比数列 \\
$\frac{1}{1-x^k}$ & $\sum_{k|n}x^n$ & 倍数 \\
$\frac{1}{(1-x)^k}$ & $\sum_{n\ge 0}\binom{n+k-1}{n}x^n$ &  \\
$\frac{x^k}{(1-x)^{k+1}}$ & $\sum_{n\ge 0}\binom{n}{k}x^n$ &  \\
$\frac{1}{\sqrt{1-4x}}$ & $\sum_{n\ge 0}\binom{2n}{n}x^n$ &  \\
$\frac{1-\sqrt{1-4x}}{2x}$ & $\sum_{n\ge 0}\binom{2n+1}{n}\frac{1}{2n+1}x^n$ & 卡特兰数\\
\bottomrule
\end{tabular}
\end{center}

\end{document}